\documentclass[12pt]{article}
\usepackage[utf8]{inputenc}
\usepackage[T1]{fontenc}
\usepackage{hyperref}
\title{Maridjabin Language Resources for K-12 Curriculum Development}
\author{Compiled by Claude Sonnet 5.5 (claude-sonnet-5-5) for Andr\'as Kornai}
\date{}
\begin{document}
\maketitle

\section{Introduction}
Maridjabin (ISO 639-3 zmj, Glottocode mari1421) is a Western Daly language of the Northern Territory, Australia, with very few elderly speakers. Sources give variant names such as Maretyabin, Maridyerbin and Marrityabin. Wikipedia treats it as a dialect of Marrithiyel. Very little material exists. No school programme or published teaching material was found in the searches.

\section{Primary Educational Resources}
No primers, textbooks or school programmes for Maridjabin were found in search results.

\section{Linguistic and Reference Materials}
\begin{itemize}
\item Wikipedia entry on Maridjabin (Marri Tjevin), with a reference to Ian Green's 1989 ANU thesis on Marrithiyel: \url{https://www.Wikipedia.org/wiki/Maridjabin}
\item Endangered Languages Project bibliography for Maridjabin: \url{https://endangeredlanguages.com/lang/1908/bibliography}
\item TLCMap gazetteer layer mentioning the language: \url{https://tlcmap.org/layers/972}
\end{itemize}

\section{Religious and Community Resources}
No religious or community-organisation resources specific to Maridjabin were found.

\section{Digital and Audio Resources}
\begin{itemize}
\item Wikitongues page for zmj: \url{https://wikitongues.org/languages/zmj/}
\end{itemize}
Search results also say that PARADISEC holds Maridjabin wordlist recordings from Ian Green's collection (1982 and 1993, Wadeye), but no direct catalogue URL appeared in the results. Access is subject to PARADISEC conditions.

\section{Curriculum Development Context}
With so few speakers and almost no published material, any curriculum would depend on archival recordings and field notes. These would need access permission and consultation with the Daly region community and language custodians. Related Marrithiyel material may help, but it should be kept distinct from Maridjabin.
\end{document}